%%%%%%%%%%%%%%%%%%%%%%%%%%%%%%%%%%%%%%%%%%%%%%%%%%%%%%%%%%%%%%%%%%%%%%%%%%%%%%
%
% 04_base_complex.tex
%
%%%%%%%%%%%%%%%%%%%%%%%%%%%%%%%%%%%%%%%%%%%%%%%%%%%%%%%%%%%%%%%%%%%%%%%%%%%%%%

\begin{mosbox}{3. The Base Complex}

\BodyFont

The foundational object of MOS is a finite simplicial complex

\[
C=(V,\Sigma),
\]

where

\[
V=\{\text{vertices}\}
\]

represents elementary concepts, while

\[
\Sigma
\]

is a family of simplices describing admissible higher-order conceptual
relationships.

Unlike an ordinary graph, simplicial complexes naturally encode
interactions involving three or more concepts simultaneously.
This richer combinatorial structure provides the geometric substrate upon
which the remainder of the MOS framework is constructed.

\end{mosbox}

%%%%%%%%%%%%%%%%%%%%%%%%%%%%%%%%%%%%%%%%%%%%%%%%%%%%%%%%%%%%%%%%%%%%%%%%%%%%%%
%% SIMPLEX HIERARCHY
%%%%%%%%%%%%%%%%%%%%%%%%%%%%%%%%%%%%%%%%%%%%%%%%%%%%%%%%%%%%%%%%%%%%%%%%%%%%%%

\vspace{4mm}

\begin{center}

\begin{tikzpicture}[
  scale=2,
  node/.style={circle, fill=MOSRed, draw=MOSGold, line width=1pt, minimum size=10pt, inner sep=0pt},
  edge/.style={draw=MOSRed, line width=2pt},
  face/.style={fill=MOSRed, opacity=0.3}
]

% Nodes
\node[node] (n1) at (0, 1.5) {};
\node[node] (n2) at (-1, 0.5) {};
\node[node] (n3) at (1, 0.5) {};
\node[node] (n4) at (-1.5, -0.5) {};
\node[node] (n5) at (0, -0.5) {};
\node[node] (n6) at (1.5, -0.5) {};
\node[node] (n7) at (0, -1.5) {};

% 2-Simplices (Faces)
\fill[face] (n2.center) -- (n4.center) -- (n5.center) -- cycle;
\fill[face] (n2.center) -- (n3.center) -- (n5.center) -- cycle;
\fill[face] (n3.center) -- (n5.center) -- (n6.center) -- cycle;
\fill[face] (n4.center) -- (n5.center) -- (n7.center) -- cycle;
\fill[face] (n5.center) -- (n6.center) -- (n7.center) -- cycle;

% 1-Simplices (Edges)
\draw[edge] (n1) -- (n2);
\draw[edge] (n1) -- (n3);
\draw[edge] (n2) -- (n3);
\draw[edge] (n2) -- (n4);
\draw[edge] (n2) -- (n5);
\draw[edge] (n3) -- (n5);
\draw[edge] (n3) -- (n6);
\draw[edge] (n4) -- (n5);
\draw[edge] (n5) -- (n6);
\draw[edge] (n4) -- (n7);
\draw[edge] (n5) -- (n7);
\draw[edge] (n6) -- (n7);

\end{tikzpicture}

\end{center}

%%%%%%%%%%%%%%%%%%%%%%%%%%%%%%%%%%%%%%%%%%%%%%%%%%%%%%%%%%%%%%%%%%%%%%%%%%%%%%
%% WHY NOT GRAPHS?
%%%%%%%%%%%%%%%%%%%%%%%%%%%%%%%%%%%%%%%%%%%%%%%%%%%%%%%%%%%%%%%%%%%%%%%%%%%%%%

\begin{highlightbox}

\SectionTitle{Graphs vs. Simplicial Complexes}

\BodyFont

\begin{center}

\begin{tabular}{p{6.5cm}|p{8.5cm}}

\textbf{Graphs}

&
\textbf{Simplicial Complexes}

\\
\hline

Pairwise interactions only

&

Interactions among arbitrary numbers of concepts

\\[3mm]

Edges represent relationships

&

Faces represent higher-order semantic structures

\\[3mm]

Limited topological information

&

Supports homology, cohomology and sheaf-theoretic constructions

\\[3mm]

Suitable for networks

&

Suitable for structured cognitive geometry

\end{tabular}

\end{center}

\end{highlightbox}

%%%%%%%%%%%%%%%%%%%%%%%%%%%%%%%%%%%%%%%%%%%%%%%%%%%%%%%%%%%%%%%%%%%%%%%%%%%%%%
%% INTUITION
%%%%%%%%%%%%%%%%%%%%%%%%%%%%%%%%%%%%%%%%%%%%%%%%%%%%%%%%%%%%%%%%%%%%%%%%%%%%%%

\begin{remarkbox}

\BodyFont

A triangle is **not** interpreted as three independent edges.

Instead, it represents a single higher-order relationship that exists only
when all three participating concepts are considered simultaneously.

This distinction is fundamental: MOS models cognition through
higher-dimensional combinatorial structures rather than solely through
pairwise associations.

\end{remarkbox}

%%%%%%%%%%%%%%%%%%%%%%%%%%%%%%%%%%%%%%%%%%%%%%%%%%%%%%%%%%%%%%%%%%%%%%%%%%%%%%
%% TRANSITION
%%%%%%%%%%%%%%%%%%%%%%%%%%%%%%%%%%%%%%%%%%%%%%%%%%%%%%%%%%%%%%%%%%%%%%%%%%%%%%

\begin{eqbox}

\[
\boxed{
\text{Geometry}
\Longrightarrow
\text{Conceptual Organization}
}
\]

\vspace{2mm}

\BodyFont

Once the combinatorial structure has been established, local mathematical
information can be attached to each simplex using a cellular sheaf. This
construction provides the second component of the MOS state space.

\end{eqbox}